%credits K. Κουνής, I. Στεφανίδου 
%refactoring : K.Draziotis (drazioti@gmail.com
%Licence CC_BY_SA

\documentclass[12pt]{article}
\usepackage{helvet} 
\usepackage{amsmath}
\usepackage{titlesec}
\usepackage{lipsum}
\usepackage{textcomp}
\usepackage{graphicx}
\graphicspath{{images/}}


\usepackage{listings}
\usepackage{xcolor}

\definecolor{codegreen}{rgb}{0,0.6,0}
\definecolor{codegray}{rgb}{0.5,0.5,0.5}
\definecolor{codepurple}{rgb}{0.58,0,0.82}
\definecolor{backcolour}{rgb}{0.95,0.95,0.92}

\lstdefinestyle{mystyle}{
    backgroundcolor=\color{backcolour},   
    commentstyle=\color{codegreen},
    keywordstyle=\color{magenta},
    numberstyle=\tiny\color{codegray},
    stringstyle=\color{codepurple},
    basicstyle=\ttfamily\footnotesize,
    breakatwhitespace=false,         
    breaklines=true,                 
    captionpos=b,                    
    keepspaces=true,                 
    numbers=left,                    
    numbersep=5pt,                  
    showspaces=false,                
    showstringspaces=false,
    showtabs=false,                  
    tabsize=2
}
\lstset{style=mystyle}

\titleformat{\section}[display]
          {\clearpage\vspace*{50pt}%
          \normalfont\huge\bfseries}%
          {{\Kappa}E{\Phi}A{\Lambda}AIO \thesection}%
          {20pt}%
          {\Huge}%
          [\vspace{40pt}]

\usepackage[algosection,commentsnumbered,ruled,vlined]{algorithm2e}
\NoCaptionOfAlgo
\usepackage{chemarrow}
\newcommand\aug{\fboxsep=-\fboxrule\!\!\!\fbox{\strut}\!\!\!}
\usepackage{graphicx}
\usepackage{gfsdidot}
\usepackage[LGR,T1]{fontenc}
\usepackage[utf8]{inputenc}
\usepackage[english,greek]{babel} % και για τις δυο γλώσσες
\usepackage{alphabeta}
\usepackage[hidelinks]{hyperref}
\usepackage{hyperref}
\usepackage{makeidx}
\usepackage{enumerate}
\usepackage{enumitem}
\usepackage{systeme}
\usepackage{algorithmic}
\usepackage{comment}
\usepackage{indentfirst}

% TEXT FORMATTING
% set spacing between lines (διάστιχο)
\usepackage{setspace}
\setstretch{1.5}
% package to customize chapters, sections and subsections style
\usepackage{titlesec}
% chapter title appearance format
\titleformat{\chapter}[display]
{\bfseries\huge}{\chaptertitlename\space\thechapter}{16pt}{}
% https://www.sharelatex.com/learn/Sections_and_chapters
\titlespacing{\chapter}{0pc}{1.5ex plus .1ex minus .2ex}{5pc}
% section title appearance format
\titleformat{\section}
{\bfseries\large}{\thesection}{14pt}{}
% subsection title appearance format
\titleformat{\subsection}
{\bfseries\normalsize}{\thesubsection}{12pt}{}
% set margins
\usepackage{geometry}
\geometry{left=3cm, right=2cm, top=4cm, bottom=3cm}
\usepackage{graphicx}
% put images in images path
\graphicspath{{images/}}
\usepackage{setspace}

% Caption customization
% use this package to set appearance for captions
\usepackage{caption}
% caption size for figures 10pt
\captionsetup[figure]{font=footnotesize,labelfont=footnotesize}
% caption size for tables 10pt and underlined
\usepackage[normalem]{ulem} % Package for underlining
\DeclareCaptionLabelFormat{label_format}{\uline{#1~#2}} % underline label
\DeclareCaptionTextFormat{text_format}{\uline{#1}} % underline text
\DeclareCaptionLabelSeparator{separator_format}{\uline{:~}} % underline separator
\captionsetup[table]{font=normalsize,labelfont=normalsize,labelformat=label_format,textformat=text_format,labelseparator=separator_format}

% use this package to define custom colors
\usepackage{xcolor}

% create colors
\colorlet{punct}{red!60!black}
\definecolor{background}{HTML}{EEEEEE}
\definecolor{delim}{RGB}{20,105,176}
\colorlet{numb}{magenta!60!black}


\usepackage{amsfonts}
\usepackage{amscd}
\usepackage{amssymb}
\newtheorem{algor}{\bf{Algorithm}}[subsection]


\newtheorem{remark}{Remark}[section]

\newtheorem{theorem}{Theorem}[section]
\newtheorem{lemma}[theorem]{Lemma}
%\newtheorem{corollary}[theorem]{Corollary}
\newtheorem{definition}[theorem]{Definition}
\newtheorem{proposition}[theorem]{Proposition}
%\theoremstyle{remark}
%\newtheorem{example}[theorem]{Example}
%\newtheorem{remark}[theorem]{Remark}
\numberwithin{equation}{section}

% use this package to show actual code listings
\usepackage{listings}

% change listings name in caption to Απεικόνιση
\renewcommand{\lstlistingname}{Απεικόνιση}

% change listings name in contents page to Κατάλογος απεικονήσεων
\renewcommand\lstlistlistingname{Κατάλογος απεικονίσεων}

% command for congruence/equivalence modulo
\newcommand{\congruence}[3]{\ensuremath{{#1}\equiv {#2}\left\pmod*{#3}\right}}
\newcommand{\congr}[2]{\ensuremath{{#1}\left\pmod*{#2}\right}\xspace}
\makeatletter
\let\@@pmod\pmod
\DeclareRobustCommand{\pmod}{\@ifstar\@pmods\@@pmod}
\def\@pmods#1{\mkern4mu({\operator@font mod}\mkern 6mu#1)}
\makeatother

% set custom colorscheme for listings with language=lang1 to make them stand out more
% http://tex.stackexchange.com/questions/83085/how-to-improve-listings-display-of-json-files
\lstdefinelanguage{lang1}{
    basicstyle=\normalfont\ttfamily,
%    numbers=left,
%    numberstyle=\scriptsize,
%    stepnumber=1,
%    numbersep=8pt,
    showstringspaces=false,
    breaklines=true,
    frame=lines,
    backgroundcolor=\color{background},
    literate=
     *{0}{{{\color{numb}0}}}{1}
      {1}{{{\color{numb}1}}}{1}
      {2}{{{\color{numb}2}}}{1}
      {3}{{{\color{numb}3}}}{1}
      {4}{{{\color{numb}4}}}{1}
      {5}{{{\color{numb}5}}}{1}
      {6}{{{\color{numb}6}}}{1}
      {7}{{{\color{numb}7}}}{1}
      {8}{{{\color{numb}8}}}{1}
      {9}{{{\color{numb}9}}}{1}
      {:}{{{\color{punct}{:}}}}{1}
      {,}{{{\color{punct}{,}}}}{1}
      {\{}{{{\color{delim}{\{}}}}{1}
      {\}}{{{\color{delim}{\}}}}}{1}
      {[}{{{\color{delim}{[}}}}{1}
      {]}{{{\color{delim}{]}}}}{1},
}

% create command for blank page
\usepackage{afterpage}
\newcommand\blankpage{%
    \null
    \thispagestyle{empty}%
    \addtocounter{page}{-1}%
    \newpage}
    
\definecolor{maroon}{HTML}{AF3235}
% add clickable hyperlinks
\usepackage{hyperref}
\hypersetup{
    colorlinks,
    citecolor=black,
    filecolor=black,
    linkcolor=black,
    urlcolor=black
}

% use fancy header and footer
\usepackage{fancyhdr}
\usepackage{blindtext} % to quickly get a full document

% Turn on the style
\pagestyle{fancy}

% Clear the header and footer
\fancyhf{}

% Set the right side of the footer to be the page number
\fancyfoot[R]{\thepage}

% set page number appearance to bottom right
\fancypagestyle{plain}{%
    \renewcommand{\headrulewidth}{0pt}
    \fancyhf{}
    \fancyfoot[R]{\thepage}%
}

\newcommand{\HRule}{\rule{\linewidth}{0.5mm}}
\newcommand{\lt}{\latintext}

\begin{document}
\begin{titlepage}

{\LARGE Αριστοτέλειο Πανεπιστήμιο Θεσσαλονίκης}
\begin{center} {\Large Σχολή Θετικών Επιστημών} \end{center}
\begin{figure}[h]
\raggedright
\hspace{90pt}
\includegraphics[width=0.5\linewidth]{logo.png}
\end{figure}
\begin{center}


\begin{figure}[h]
\centering 
%\includegraphics[width=0.3\linewidth]{authLogo.png}
\end{figure}
\begin{center}
% leave 2 cm from above text

\HRule \\[0.4cm]
{\huge Υλοποίηση \lt{Quadratic Sieve}}

\HRule \\[0.4cm]
\end{center}

% put this on the bottom
\vfill
\begin{doublespacing}

{\LARGE 
Καλαθάς Μάνος - ΑΕΜ:101 \par}

\vfill 
{\Large \today}
\end{doublespacing}
\end{center}
\end{titlepage}

% insert table of contents
\tableofcontents

\clearpage

% leave blank page before main part
\blankpage

\section {\lt{Quadratic Sieve}}

\subsection{Περιγραφή προγράμματος}
Το πρόγραμμα δέχεται ως είσοδο από τον χρήστη τον αριθμό που επιθυμεί να παραγοντοποιήσει και ως μεταβλητές έχει ένα άνω όριο B και το μήκος του \lt{sieve}. Ο χρήστης μπορεί να αλλάξει αυτές τις τιμές όπως επιθυμεί και κάθε φορά που ο αλγόριθμος αποτυγχάνει να βρει μία λύση διπλασιάζει το όριο B.

\subsection {Αρχικοποίηση}
Στην αρχικοποίηση του αλγορίθμου υπολογίζονται όλοι οι πρώτοι αριθμοί μικρότεροι του B για τους οποίους ισχύει η σχέση legendre \(\frac{n}{p} = 1\).

\subsection{Κοσκίνισμα}
Βάση των πρώτων που βρήκαμε στο προηγούμενο βήμα, πραγματοποιείται τώρα ένα κοσκίνισμα των αριθμών \((x^2 -n)\) με το x να παίρνει τιμές από την ρίζα του αριθμού που θέλουμε να παραγοντοποιήσουμε μέχρι το διάστημα που ορίζει η μεταβλητή στην αρχή του προγράμματος. Οι αριθμοί που βρίσκουμε σε αυτό το στάδιο λέγονται B-ομαλοί και έχουν την ιδιότητα να μην έχουν πρώτους παράγοντες μεγαλύτερους του B. Το πλήθος τους είναι όσο το μέγεθος της βάσης των πρώτων που υπολογίστηκαν νωρίτερα, συν ένα.

\subsection{Γραμμική Άλγεβρα}
Αρχικά φτιάχνουμε έναν πίνακα όπου κάθε γραμμή είναι ένα διάνυσμα των εκθετών όλων των πρώτων της βάσης που βρήκαμε modulo 2. Στην συνέχεια φέρνουμε τον πίνακα σε reduced echelon form με την χρήση απαλοιφής του Gauss στο πεδίο GF(2) και βρίσκουμε ποιες γραμμές είναι γραμμικά εξαρτημένες μεταξύ τους.

\subsection{Παραγοντοποίηση}
Έχοντάς βρει ένα συνδυασμό εξαρτημένων γραμμών από το προηγούμενο βήμα, και έστω \((x_1,x_2,..,x_k)\) αυτές παίρνουμε την τιμή \(x = x_1,x_2,..,x_k \pmod{n}\) και \(y = \sqrt{(x_1^2-n)(x_2^2-n)...(x_k^2-n)} \pmod{n}.\) τις οποίες έχουμε υπολογίσει και νωρίτερα στον αλγόριθμο κατά την εύρεση B ομαλών στοιχείων.
Βρίσκουμε το gcd(x-y, n) και αν ο μέγιστος κοινός διαιρέτης έχει τιμή διάφορη από το n και το 1, έχουμε βρει μία παραγοντοποίηση του n.

\subsection{Σχόλια}
Δεν κατάφερα να κάνω τον αλγόριθμο να παραγοντοποιεί RSA 100 αριθμούς, αποφάσισα παρ' όλα αυτά να σας τον στείλω.

\bibliographystyle{plain}
\bibliography{bib.bib}

\end{document}